\captionof{figure}{\textbf{Low functional and anatomical co-correspondence except for a strong cortical--subcortical division.}
\textbf{a} Task-aligned response vectors for 53,021 neurons, ordered by Rastermap and overlaid with Beryl anatomical colors.
\textbf{b} The same response vectors ordered by anatomical region, illustrating that functional response motifs recur across regions.
\textbf{c} Response vectors from six example regions (CP, caudoputamen; MRN, midbrain reticular nucleus; ZI, zona incerta; MOp, primary motor area; CA1, hippocampal field CA1; CUL4,5, cerebellar culmen lobules IV--V), retaining their Rastermap order and colored by position in that ordering.
\textbf{d} UMAP embedding of neural response vectors colored by Beryl region, with anatomical coordinates shown alongside.
\textbf{e} UMAP embedding and anatomical coordinates colored by 25 functional k-means clusters.
\textbf{f} Beryl-region composition of each functional cluster.
\textbf{g} Functional-cluster frequencies for eight example Beryl regions (PA, PAA, MOB, MEA, MRN, SCm, PRNr, PGRN). Each inset gives its regional specialization score, $S=1-H(p)/\log(25)$, where $p$ is the region's distribution across clusters and $H$ is Shannon entropy.
\textbf{h} Distributions of regional specialization estimated from functional clustering and from decoding.
\textbf{i} Regional specialization and coarse anatomical partitions displayed on Swanson flatmaps.
\textbf{j} Clustering-based specialization correlates strongly with decoding-based specialization across regions (log--log axes; the eight regions of g are labelled).}
\label{fig:anatomy_function}
